\documentclass[10pt,a4paper]{amsart}
\usepackage[margin=19mm]{geometry}
\usepackage{amsmath,amssymb,mathtools}
\usepackage[scr=rsfs,cal=euler]{mathalfa}
\usepackage{xfrac}
\usepackage[unicode,hidelinks,bookmarks=false]{hyperref}
\newcommand{\ie}{i.e.\ }
\newcommand{\cf}{cf.\ }
\newcommand{\resp}{respectively\ }
\newcommand{\Subsection}{\subsection}
\providecommand{\nobreakdash}{\nobreak\hbox{-}}
\setlength{\emergencystretch}{2em}
\multlinegap0pt
\usepackage{bm}
\usepackage{bbm}
\usepackage{dsfont}
\usepackage{xspace}
\usepackage{cancel}
\usepackage{mathtools}
\usepackage{amsthm}
\makeatletter
\@ifundefined{prop}{\newtheorem{prop}{Proposition}}{}
\makeatother
\title[An Elliptic Curve Governing Hopf Linking in an $A\_4$-Symmetr]{An Elliptic Curve Governing Hopf Linking in an $A\_4$-Symmetric Tensegrity}
\author{Taizo Sadahiro}
\date{Source: arXiv:2604.18116v1 (CC BY 4.0)}
\begin{document}
\maketitle

\noindent\emph{Source and licence.} An Elliptic Curve Governing Hopf Linking in an $A_4$-Symmetric Tensegrity. Version arXiv:2604.18116v1 (CC BY 4.0), distributed under the Creative Commons Attribution 4.0 International licence, as recorded in the arXiv licence metadata for this exact version and shown on the abstract page https://arxiv.org/abs/2604.18116v1. Full rights notice: licenses/arxiv-2604.18116-CC-BY-4.0.txt.\par\smallskip

\noindent\textit{Editorial scope.} Counted unit: the Prop whose source label is \texttt{prop:onehopflink} in the article (source0.tex lines 380--393, source label \texttt{prop:onehopflink}), delivered together with the article's own complete original proof of it (source0.tex lines 394--509, one \texttt{proof} environment of 116 lines). No mathematics is written, repaired or re-derived: statement and proof are the article's own text line for line. Required context retained uncounted: none. Every definition, display and label of those lines is retained unchanged. Omitted from this package and not redistributed: 1--379, 510--1096. Nothing inside a retained excerpt is omitted or altered: every retained body is delivered exactly as the article's lines are written, so no omission receipt of any kind accompanies this package. Citations: the retained excerpts carry no citation command at all, so no bibliography entry of the article is transcribed and no reference is redistributed. Reference adaptation: none was needed and none was made; every reference inside the delivered unit resolves inside it, so no unresolved reference or citation is produced. Printed slips kept literal: no printed word, symbol, hypothesis or number of the counted statement or of its proof was altered. Layout and class: the article's own class is \texttt{article} with packages including amsmath, amssymb, amsthm, float, url, tikz, pgf, tikz-3dplot, graphicx; the wrapper is the lane's pinned \texttt{amsart} editorial wrapper at 10pt on A4, which restates the theorem-like environments the retained text needs and adds the article's own macro definitions that the retained text needs (none), with extra wrapper packages (bm, bbm, dsfont, xspace, cancel, mathtools). The delivered numbering is the unit's own. Assets: no retained span contains a figure environment, a graphics-inclusion command, a PDF-page inclusion, a table or a TikZ picture; no image file is copied into this package, the package declares no asset dependency and no image file is redistributed with it. Licence: version arXiv:2604.18116v1 of this article is distributed under the Creative Commons Attribution 4.0 International licence, as recorded in the arXiv licence metadata for this exact version and shown on the abstract page https://arxiv.org/abs/2604.18116v1. A full rights notice is delivered at licenses/arxiv-2604.18116-CC-BY-4.0.txt. Characters: every retained excerpt is pure ASCII, so the delivered engine, XeLaTeX with its default Latin Modern text font, reports no missing character.
\begin{prop}
\label{prop:onehopflink}
At the distinguished point
$
(x,y)=\left(\frac12,-\frac13\right),
$
the two triangular strut components
\[
\Delta=\{{\mathbf p}_0,\rho(s){\mathbf p}_0,\rho(s^2){\mathbf p}_0\}
\quad\text{and}\quad
\Delta'=\{\rho(c_1){\mathbf p}_0,\rho(c_1s){\mathbf p}_0,\rho(c_1s^2){\mathbf p}_0\}
\]
form a Hopf link.
\end{prop}

\begin{proof}
At the distinguished point \((x,y)=\left(\frac12,-\frac13\right)\), the general expression for the null vector gives
\[
{\mathbf p}_0=\left(0,\frac53,\frac53\right).
\]
Since the realization is unchanged up to an overall scalar multiple, we normalize and write
\[
{\mathbf p}_0=(0,1,1).
\]
Let
\[
A={\mathbf p}_0=
\begin{pmatrix}
0\\ 1\\ 1
\end{pmatrix},\qquad
B=\rho(s){\mathbf p}_0=
\begin{pmatrix}
1\\ -1\\ 0
\end{pmatrix},\qquad
C=\rho(s^2){\mathbf p}_0=
\begin{pmatrix}
-1\\ 0\\ -1
\end{pmatrix},
\]
so that
\[
\Delta=\triangle ABC.
\]
Similarly, let
\[
D=\rho(c_1){\mathbf p}_0=
\begin{pmatrix}
-1\\ 1\\ 0
\end{pmatrix},\qquad
E=\rho(c_1s){\mathbf p}_0=
\begin{pmatrix}
1\\ 0\\ -1
\end{pmatrix},\qquad
F=\rho(c_1s^2){\mathbf p}_0=
\begin{pmatrix}
0\\ -1\\ 1
\end{pmatrix},
\]
so that
\[
\Delta'=\triangle DEF.
\]

The plane containing \(\Delta\) is
\[
\Pi:\ X+Y-Z=0.
\]
Indeed, the three vertices \(A,B,C\) satisfy
\[
0+1-1=0,\qquad 1-1-0=0,\qquad -1+0-(-1)=0.
\]

We claim that the edge \(EF\) meets the interior of the triangular disk bounded by \(\Delta\).
Its midpoint is
\[
\frac12E+\frac12F
=
\frac12
\begin{pmatrix}
1\\ -1\\ 0
\end{pmatrix}
=
\begin{pmatrix}
\frac12\\ -\frac12\\ 0
\end{pmatrix},
\]
and this point lies on \(\Pi\), since
\[
\frac12-\frac12-0=0.
\]
Moreover,
\[
\frac12E+\frac12F
=
B+\frac16(A-B)+\frac16(C-B).
\]
Since
\[
\frac16>0,\qquad
1-\frac16-\frac16=\frac23>0,
\]
this point is a strict convex combination of \(A,B,C\). Hence it lies in the interior of the triangular disk bounded by \(\Delta\). Therefore the edge \(EF\) intersects the interior of that disk.
In particular, the intersection point lies on the symmetry axis of the equilateral triangle 
$\Delta$ through the vertex B.

Next we show that the other two edges of \(\Delta'\) do not meet the triangular disk bounded by \(\Delta\).
First,
\[
D = B + \frac23(A-B) +\frac23(C-B),
%A-(B-A)+\frac23(C-A),
\]
so \(D\in \Pi\), but \(D\) lies outside \(\Delta\), because one barycentric coefficient is negative.

Now consider the function
\[
\ell(X,Y,Z)=X+Y-Z.
\]
Since \(\Pi\) is given by \(\ell=0\), we have
\[
\ell(D)= -1+1-0=0,\qquad
\ell(E)=1+0-(-1)=2,\qquad
\ell(F)=0+(-1)-1=-2.
\]
Thus the segment \(DE\) meets the plane \(\Pi\) only at the endpoint \(D\), and likewise the segment \(DF\) meets \(\Pi\) only at the endpoint \(D\). Since \(D\) lies outside \(\Delta\), neither \(DE\) nor \(DF\) intersects the triangular disk bounded by \(\Delta\).

Consequently, the closed polygon \(\Delta'\) meets the triangular disk bounded by \(\Delta\) in exactly one point, namely the interior point of the edge \(EF\) described above. Therefore
\[
\bigl|\operatorname{lk}(\Delta,\Delta')\bigr|=1.
\]
Since both \(\Delta\) and \(\Delta'\) are embedded triangles, each is an unknot. Hence \(\Delta\) and \(\Delta'\) form a Hopf link.
\end{proof}

\end{document}
